% ============================================================
% 20_tables.tex — Tabellen-Helfer (semantisches Table-System)
% ============================================================
% Zentrales Design-Prinzip:
%   1. Tabellen fuellen IMMER \linewidth (kein manuelles p{0.16..}).
%   2. Spaltenbreiten werden PROPORTIONAL deklariert (Y{0.18} L{0.54} ...),
%      sodass die Summe der Faktoren = Anzahl X-Spalten ergibt.
%   3. Zebra-Streifen + Header-Highlight sind ins Environment eingebaut;
%      direkte Farb-Primitive sind in Kapiteln per Linter verboten.
%
% Damit ist strukturell ausgeschlossen:
%   - zu kurze Zebra-Streifen (Tabelle schmaler als Box),
%   - überstehende Zeilen (Content sprengt Box-Breite),
%   - Drift zwischen Kapiteln (jede Tabelle nutzt dieselben Wrapper).

% ------------------------------------------------------------
% Spaltentypen (tabularx-X-Varianten)
% ------------------------------------------------------------
% Feste Varianten — jede zaehlt als Faktor 1:
% m{#1}: X-Spalten vertikal mittig (Text + Grafik in einer Zeile).
\renewcommand{\tabularxcolumn}[1]{m{#1}}
\newcolumntype{L}{>{\RaggedRight\arraybackslash}X}   % linksbuendig, umbricht
\newcolumntype{C}{>{\Centering\arraybackslash}X}     % zentriert,  umbricht
\newcolumntype{R}{>{\RaggedLeft\arraybackslash}X}    % rechtsbuendig, umbricht

% Parametrisierte, proportionale Varianten — #1 = relativer Anteil.
%
% WICHTIG (tabularx-Konvention): Die Summe aller Faktoren in EINER
% Tabelle muss = Anzahl X-Spalten sein, damit die Tabelle \linewidth
% voll ausfuellt. L/C/R zählen als Faktor 1.
%
% Beispiele (2 Spalten → Summe 2):
%   Y{0.6}  Y{1.4}            → 30% / 70%
%   Y{1}    Y{1}              → 50% / 50% (== L L)
%   Q{0.8}  Y{1.2}            → 40% / 60%, Spalte 1 rechtsbuendig
% Beispiele (3 Spalten → Summe 3):
%   Y{0.6}  Y{1.5}  Y{0.9}    → 20% / 50% / 30%
%   L       Y{2}    L         → 25% / 50% / 25%
% Beispiele (4 Spalten → Summe 4):
%   Y{1.2}  Y{0.8}  Y{1.2}  Y{0.8}  → 30% / 20% / 30% / 20%
\newcolumntype{Y}[1]{>{\hsize=#1\hsize\RaggedRight\arraybackslash}X} % linksb., proport.
\newcolumntype{Z}[1]{>{\hsize=#1\hsize\Centering\arraybackslash}X}   % zentriert, proport.
\newcolumntype{Q}[1]{>{\hsize=#1\hsize\RaggedLeft\arraybackslash}X}  % rechtsb., proport.

% Zentrierte Formel-Spalte (Math-Mode automatisch, proportional)
\newcolumntype{F}[1]{>{\hsize=#1\hsize\Centering\arraybackslash$}X<{$}}

% ------------------------------------------------------------
% Zebra-Streifen (Scannbarkeit)
% ------------------------------------------------------------
\newcommand{\ZSFzebraBG}{\TableZebraColor}
\newcommand{\SetZSFzebraBG}[1]{\renewcommand{\ZSFzebraBG}{#1}}
% Eine Startzeile-Variante genügt; die Umgebung leitet die Zeile aus der
% Kopfzeilen-Option ab. Ein zusätzliches \ZSFzebra mit fest verdrahteter 2
% wäre ein zweiter Weg zum selben Ergebnis (rules/02_ai_mandate).
\newcommand{\ZSFzebraStart}[1]{\rowcolors{#1}{\BoxPlainBackColor}{\ZSFzebraBG}}
\newcommand{\ZSFzebraOff}{\rowcolors{1}{}{}}

% ------------------------------------------------------------
% Header-Highlight
% ------------------------------------------------------------
% Konvention:
%   \ZSFheaderRow   — startet die Header-Zeile und setzt den Zeilen-Hintergrund.
%   \ZSFhead{<Text>} — umschliesst den Inhalt JEDER Header-Zelle (Farbe + bold).
%
% Warum zweigeteilt? In tabularx-X-Spalten propagiert \color, das direkt nach
% \rowcolor am Zeilenanfang gesetzt wird, NICHT zuverlässig über die
% &-Grenzen. Die erste Zelle verliert die Farbe, andere behalten sie — je
% nach Spaltenbreite und Content (siehe SI-Praefix-Bug). Per-Zelle-Wrapper
% via \ZSFhead{...} ist robust und konsistent.
\newcommand{\ZSFheaderRowBG}{\chaptercolor}
\newcommand{\ZSFheaderRowFG}{\chaptercolortext}
\newcommand{\ZSFheaderRow}{\rowcolor{\ZSFheaderRowBG}}
% \ZSFInkOwned steht direkt neben dem \color, das den Kontrast setzt: Die
% Kopfzelle ist eine gesaettigte Flaeche mit eigener Textfarbe, eine
% Inline-Farbe darin waere unlesbar (40_colors_structure -> Ink-Vertrag).
\newcommand{\ZSFhead}[1]{{\color{\ZSFheaderRowFG}\ZSFInkOwned\ZSFfontTableHead #1}}

% ------------------------------------------------------------
% Zellen verbinden
% ------------------------------------------------------------
% \ZSFspan{<L|C|R>}{<Anzahl Spalten>}{<Inhalt>}
%
% Eine Zelle über mehrere Spalten — die Ebene, die zwischen Tabelle und Zelle
% fehlte. Drei Rollen, ein Baustein, weil es technisch dieselbe Sache ist und
% nur der Inhalt sie unterscheidet (rules/02_ai_mandate → Katalog-Ökonomie):
%
%   gruppierter Kopf      \ZSFheaderRow \ZSFhead{Querschnitt} &
%                           \ZSFspan{C}{2}{\ZSFhead{Biegung}} & …
%   volle Breite          \ZSFspan{C}{3}{$\displaystyle \cos(2x)=\dots$} \\
%   Abschnittszeile       \ZSFspan{L}{2}{\ZSFlabel{Rang 2: voller Rang}} \\
%
% Ersetzt rohes \multicolumn, das in Kapiteln verboten ist und für das es bis
% hierher keinen Ersatz gab — die einzige Stelle im System, an der eine
% belegte Notwendigkeit gar keinen legalen Weg hatte.
%
% WARUM die Ausrichtung ein PFLICHTargument ist und kein optionaler Schlüssel
% wie sonst überall: TeX entscheidet am Zellenanfang, ob die Spaltenvorlage
% eingesetzt wird, und lässt sich davon nur durch ein \omit abbringen — das
% \multicolumn als erstes setzt. Auf dem Weg dorthin darf TeX ausschliesslich
% EXPANDIEREN. Jede Form mit optionalem Argument führt aber eine Zuweisung aus,
% bevor \multicolumn erreicht ist (xparse ebenso wie \newcommand[..][..]), und
% die Vorlage steht dann schon: "Misplaced \omit". Gemessen an beiden Formen,
% nicht hergeleitet. Deshalb ist der Weg hierher bewusst der einzige im System,
% der von der Schlüsselkonvention abweicht — mit sichtbarem Grund, damit ihn
% niemand später "geradezieht" und den Fehler zurückholt.
%
% Dieselbe Einschränkung erzwingt die Form darunter: Der Spaltenbuchstabe muss
% LITERAL bei \multicolumn ankommen. Ein Makro, das ihn liefert, geht nicht —
% der Preamble-Parser von `array` vergleicht Tokens und expandiert dabei nicht
% ("Illegal pream-token"). Auch das gemessen. Deshalb steht die Auswahl als
% \if-Kette VOR \multicolumn: Sie ist rein expandierbar, hält die \omit-Regel
% ein und übergibt am Ende ein echtes l/c/r.
%
% Der Fehlerfall läuft absichtlich in den Zellinhalt und nicht davor: Dort ist
% \omit schon gesetzt, eine Meldung ist damit unschädlich. Ein stiller Rückfall
% auf `C` sähe im PDF wie eine bewusst zentrierte Zelle aus.
%
% Eine verbundene Zelle bricht NICHT um: l/c/r statt einer X-Spalte, weil die
% Breite eines Verbunds erst beim Satz feststeht. Das deckt die Rollen oben ab
% (Kopf, Formel, kurze Abschnittszeile); für umbrechenden Fliesstext über die
% volle Breite ist die Tabelle der falsche Baustein.
\makeatletter
\newcommand{\ZSF@spanBadAlign}[1]{%
  \PackageError{ZSF}{Unbekannte Ausrichtung '#1' in \string\ZSFspan}%
    {Gueltig sind L, C und R -- dieselben Buchstaben wie bei den
     Spaltentypen L/C/R.}%
}
\newcommand{\ZSFspan}[3]{%
  \if L#1\multicolumn{#2}{l}{#3}%
  \else\if R#1\multicolumn{#2}{r}{#3}%
  \else\if C#1\multicolumn{#2}{c}{#3}%
  \else\multicolumn{#2}{c}{\ZSF@spanBadAlign{#1}#3}%
  \fi\fi\fi
}
\makeatother

% ------------------------------------------------------------
% Semantische Tabellen-Umgebung (Single Point of Truth)
% ------------------------------------------------------------
% EINE Umgebung mit Reglern statt mehrerer fast gleicher Varianten
% (rules/02_ai_mandate → Katalog-Ökonomie): Kopfzeile und Zebra sind zwei
% Ja/Nein-Fragen, keine drei Umgebungsnamen.
%
%   \begin{ZSFtable}[<Optionen>]{<colspec>}
%
% Optionen (Komma-Liste, alles weglassbar):
%   header      Kopfzeile vorhanden (Default). 'header=false': keine.
%   zebra       Zebra-Streifen an (Default). 'zebra=false': aus.
%   font=dense  eine Stufe kleiner, für lange Tabellen mit knappem
%               Zellinhalt. Default 'font=normal'.
%   rows        Zeilenhöhe: 'normal' (Default) | 'roomy' | 'tight'.
%   colsep      Zellpolsterung: 'normal' (Default) | 'tight'.
%
% Warum `font`, `rows` und `colsep` drei Regler sind und nicht einer: Sie
% beantworten drei verschiedene Fragen, die sich kreuzen. Eine Formeltabelle
% will grosse Schrift UND viel Zeilenhöhe; ein siebenspaltiges Register will
% kleine Schrift, knappe Zellpolsterung und normale Zeilenhöhe. Ein
% gemeinsamer „Dichte"-Regler könnte nur eine dieser Kombinationen treffen und
% müsste die anderen still verwerfen (rules/21_box_options).
%
% Die Zebra-Startzeile folgt automatisch der Kopfzeile: mit Header ab
% Zeile 2, ohne ab Zeile 1.
%
% colspec nutzt L/C/R (gleichverteilt, jeweils Faktor 1) oder
% Y{.}/Z{.}/Q{.}/F{.} (proportional). Summe der Faktoren muss gleich
% der Anzahl X-Spalten sein (tabularx-Konvention).
%
% Beispiele:
%   \begin{ZSFtable}{Y{0.6} Y{1.55} Y{0.85}}            % Header + Zebra
%   \begin{ZSFtable}[header=false]{Y{1.2} Q{0.8}}       % Zebra ab Zeile 1
%   \begin{ZSFtable}[header=false, zebra=false]{L C}    % nackt
%   \begin{ZSFtable}[font=dense]{Y{0.6} Y{1.55} Y{0.85}}
\makeatletter
\newif\ifZSF@tblHeader
\newif\ifZSF@tblZebra
\newcommand{\ZSF@tblFont}{\ZSFfontContent}
\newcommand{\ZSF@tblStretch}{\ZSFtableArrayStretch}
\newlength{\ZSF@tblColSep}
\newcommand{\ZSF@tblSpec}{}
\newcommand{\ZSF@beginTabularx}[1]{\tabularx{\linewidth}{#1}}
\pgfkeys{
  /zsf/table/.is family, /zsf/table,
  header/.is if=ZSF@tblHeader,
  zebra/.is if=ZSF@tblZebra,
  font/.is choice,
  font/normal/.code={\renewcommand{\ZSF@tblFont}{\ZSFfontContent}},
  font/dense/.code={\renewcommand{\ZSF@tblFont}{\ZSFfontContentDense}},
  rows/.is choice,
  rows/normal/.code={\renewcommand{\ZSF@tblStretch}{\ZSFtableArrayStretch}},
  rows/roomy/.code ={\renewcommand{\ZSF@tblStretch}{\ZSFtableArrayStretchRoomy}},
  rows/tight/.code ={\renewcommand{\ZSF@tblStretch}{\ZSFtableArrayStretchTight}},
  % `colsep` regelt ausschliesslich den Abstand ZWISCHEN den Spalten. Der
  % Abstand zur Boxkante haengt nicht daran — siehe der @-Eintrag unten.
  colsep/.is choice,
  colsep/normal/.code={\setlength{\ZSF@tblColSep}{\ZSFtableEdgePad}},
  colsep/tight/.code ={\setlength{\ZSF@tblColSep}{\ZSFtableEdgePadTight}},
  ZSF@tblDefaults/.style={header=true, zebra=true, font=normal,
                          rows=normal, colsep=normal},
}
\NewDocumentEnvironment{ZSFtable}{O{} m}{%
  \begingroup
  \pgfkeys{/zsf/table, ZSF@tblDefaults, #1}%
  \setlength{\tabcolsep}{\ZSF@tblColSep}%
  \renewcommand{\arraystretch}{\ZSF@tblStretch}%
  \ZSF@tblFont
  \ifZSF@tblZebra
    \ifZSF@tblHeader\ZSFzebraStart{2}\else\ZSFzebraStart{1}\fi
  \else
    \ZSFzebraOff
  \fi
  % Der Abstand zur BOXKANTE, getrennt vom Abstand zwischen den Spalten.
  %
  % \tabcolsep trug beide Rollen: In einer randlosen tablebox (padx=none, damit
  % Zebra und Kopfzeile bis an den Rand laufen) ist der aeussere Rand genau ein
  % \tabcolsep. Wer die Spalten enger stellte, schob damit den Text an die
  % Kante — und das ist gerade bei der vielspaltigen Tabelle noetig, fuer die
  % `colsep=tight` existiert.
  %
  % Der @-Eintrag ersetzt den \tabcolsep an der Aussenkante durch einen eigenen
  % Wert. Bei `colsep=normal` ist er wertgleich, die Geometrie aller
  % bestehenden Tabellen bleibt damit unveraendert; bei `colsep=tight` haelt er
  % den Rand, waehrend innen gekuerzt wird. Deshalb ohne Fallunterscheidung —
  % ein \if um \tabularx herum bricht ohnehin, weil die Umgebung ihren Rumpf
  % selbst bis \end{tabularx} liest und das \fi dabei stehenbleibt.
  %
  % Die Spaltenaufteilung wird VORHER aufgeloest. Der Preamble-Parser von
  % `array` vergleicht Tokens und expandiert dabei nicht: Ein Makro als colspec
  % meldet "Illegal pream-token" (gemessen, nicht hergeleitet). Genau das
  % braucht aber die Boxgruppe, deren Mitglieder ihre gemeinsame Aufteilung als
  % \ZSFgroupcols lesen, statt sie in jedem Block zu wiederholen.
  % Der \edef ist unbedenklich: Spaltenbuchstaben sind Zeichen (L/C/R/Y…
  % stammen aus \newcolumntype und sind keine Makros), \hskip ist ein Primitiv
  % und \ZSFtableEdgePad ein Register — nichts davon expandiert.
  \edef\ZSF@tblSpec{@{\hskip\ZSFtableEdgePad}#2@{\hskip\ZSFtableEdgePad}}%
  \expandafter\ZSF@beginTabularx\expandafter{\ZSF@tblSpec}%
}{%
  \endtabularx
  \endgroup
}
\makeatother
